\addcontentsline{toc}{chapter}{Introduction}
\chapter*{Introduction}

The studied domain of this project sits at the intersection of programming language design 
and educational technology, specifically focused on the use of domain-specific languages as 
tools for teaching computational thinking to children. As digital literacy becomes an 
increasingly essential skill, the question of how to introduce programming concepts to young 
learners in an engaging and effective way has attracted significant research attention across 
both computer science and pedagogy.

The motivation for this project stems from a clear gap in the existing landscape of 
educational programming tools. Block-based environments such as Scratch successfully lower 
the entry barrier but create a problematic disconnect from real text-based programming, 
making the eventual transition to languages like Python or Java unnecessarily abrupt. On the 
other end of the spectrum, tools that operate within Minecraft --- a game already deeply 
familiar to the target age group --- such as ComputerCraft, impose technical complexity far 
beyond what a beginner can reasonably handle. CodeCraft was conceived to occupy the space 
between these two extremes: a text-based language simple enough for a ten-year-old, yet 
grounded in real programming syntax and capable of producing immediate, visible results 
inside a game children already care about.

The relevance of the topic is supported by a growing body of empirical research demonstrating 
that game-based, visually immediate programming environments significantly improve motivation 
and computational thinking outcomes among young learners. Minecraft in particular has seen 
documented adoption across primary and secondary education in over one hundred countries, yet 
accessible programmatic scripting tools tailored to novice children remain scarce, making 
this a timely and practical contribution.

The general objectives of the project are threefold: to design a domain-specific language 
suited to children aged 10--15 with no prior programming experience, to define its formal 
grammar and core language constructs, and to implement the foundational components of its 
toolchain, namely the lexer and parser.

The research methodology combined a structured literature review of educational technology, 
programming pedagogy, and DSL design with an empirical-informed design process, drawing on 
studies covering learning curves across programming paradigms, identifier readability, and 
common beginner mistakes to justify each design decision.

The report is organized into two main chapters. Chapter 1 presents the domain analysis, 
covering the educational context, a critical examination of existing tools, their 
shortcomings, the proposed solution and its use case scenarios, and a profile of the target 
users and stakeholders. Chapter 2 addresses the language design, detailing the DSL 
requirements, the formal BNF grammar, and the implementation of the lexer and parser 
components. The report concludes with a summary of outcomes and directions for future work.
